\section{Framework teorico}

\subsection{Dinamica del sottostante}

Nel modello di Black-Scholes, il prezzo del sottostante segue un moto browniano geometrico sotto misura risk-neutral:
Questa formulazione e' standard nella modellazione continua dei mercati finanziari \cite{shreve2004stochastic}.

\begin{equation}
    dS_t = r S_t dt + \sigma S_t dW_t,
\end{equation}

dove $S_t$ e' il prezzo dell'asset al tempo $t$, $r$ e' il tasso risk-free, $\sigma$ e' la volatilita' costante e $W_t$ e' un moto browniano standard.

La soluzione della SDE e':

\begin{equation}
    S_T = S_0 \exp\left[\left(r - \frac{1}{2}\sigma^2\right)T + \sigma \sqrt{T} Z\right],
    \qquad Z \sim \mathcal{N}(0,1).
\end{equation}

\subsection{Opzione call europea}

Una call europea con strike $K$ e maturity $T$ ha payoff:

\begin{equation}
    \max(S_T - K, 0).
\end{equation}

Sotto pricing risk-neutral, il prezzo e' il valore atteso scontato del payoff:

\begin{equation}
    C = e^{-rT}\mathbb{E}^{\mathbb{Q}}\left[\max(S_T - K, 0)\right].
\end{equation}

\subsection{Formula di Black-Scholes}

La formula analitica per una call europea e':

\begin{equation}
    C_{BS} = S_0 N(d_1) - K e^{-rT} N(d_2),
\end{equation}

dove $N(\cdot)$ e' la funzione di ripartizione della normale standard e:

\begin{align}
    d_1 &= \frac{\ln(S_0/K) + \left(r + \frac{1}{2}\sigma^2\right)T}{\sigma\sqrt{T}}, \\
    d_2 &= d_1 - \sigma\sqrt{T}.
\end{align}

Nel progetto, $C_{BS}$ viene usato come valore target per il training e come riferimento per la valutazione.

\subsection{Pricing Monte Carlo}

La simulazione Monte Carlo approssima il valore atteso risk-neutral generando molte realizzazioni di $S_T$:

\begin{equation}
    \hat{C}_{MC} =
    e^{-rT}
    \frac{1}{M}
    \sum_{i=1}^{M}
    \max(S_T^{(i)} - K, 0).
\end{equation}

Per $M$ sufficientemente grande, $\hat{C}_{MC}$ converge a $C_{BS}$.
Nel progetto Monte Carlo viene usato come benchmark numerico e come confronto rispetto alla rete neurale.
